% --- Archon physics formalization source begin ---
% archon:physics
% archon:covers PhyXMiniProblems/problem_phyx_mini_0559.lean
% archon:source-report reports/phyx_mini/problem_phyx_mini_0559.source.json

\chapter{Physics problem phyx\_mini\_0559}
\label{ch:PhyXMiniProblems_problem_phyx_mini_0559}

\paragraph{Problem source.}
\#\# Physical scenario

A rectangular box at rest in S' has sides \textbackslash{}( a' = 2 \textbackslash{}text\{ m\} \textbackslash{}), \textbackslash{}( b' = 2 \textbackslash{}text\{ m\} \textbackslash{}), and \textbackslash{}( c' = 4 \textbackslash{}text\{ m\} \textbackslash{}) and is oriented as shown in figure S' moves with \textbackslash{}( \textbackslash{}beta = 0.65 \textbackslash{}) with respect to the laboratory frame S.

\#\# Question

Compute the volume of the box in S' and in S.

\#\# Answer choices

- A: A: \textbackslash{}( : 18.6m\textasciicircum{}3 \textbackslash{})

- B: B: \textbackslash{}( 21.7 \textbackslash{})

- C: C: \textbackslash{}( 15.4m\textasciicircum{}3 \textbackslash{})

- D: D: \textbackslash{}( 12.2m\textasciicircum{}3 \textbackslash{})

\#\# Figure caption (auxiliary; use the image as primary evidence)

The image shows a three-dimensional geometric shape resembling a parallelepiped situated within a coordinate system. The axes are labeled \textbackslash{}(x'\textbackslash{}), \textbackslash{}(y'\textbackslash{}), and \textbackslash{}(z'\textbackslash{}), indicating a coordinate space. 

The parallelepiped is shaded in blue, with its vertices labeled \textbackslash{}(a'\textbackslash{}), \textbackslash{}(b'\textbackslash{}), and \textbackslash{}(c'\textbackslash{}). These labels likely represent vectors or points along the axes. The label \textbackslash{}(S'\textbackslash{}) appears near the top of the shape, possibly marking a specific point or surface.

A number "25" is present, possibly indicating a measurement along one axis or representing a dimension of the shape. Solid and dashed lines are used to illustrate the edges and axis relationships, indicating spatial orientation and perspective of the shape in the coordinate system.

\#\# Dataset metadata

Relativity; Spatial Relation Reasoning, Physical Model Grounding Reasoning

\paragraph{Recorded answer/context.}
D: D: \textbackslash{}( 12.2m\textasciicircum{}3 \textbackslash{})

\paragraph{Figure/image path.}
/root/proposal\_for\_physic/science-mango/phyx\_mini\_run/phyx\_data/test\_image/559.png

\paragraph{Formalization target.}
create a compiling Lean file with sorry bodies at `PhyXMiniProblems/problem\_phyx\_mini\_0559.lean`. The Lean declarations must preserve the physical quantities, dimensions or dimensional roles, figure labels, governing-law hypotheses, and final relation expressed by this problem.
Use Mathlib/Physlib names found through LeanExplore where available. If a physics API is missing, introduce faithful local abstractions rather than scalar placeholder aliases.

\begin{theorem}[Physics formalization target]
\label{thm:physics:phyx_mini_0559:target}
\lean{PhyXMiniProblems.ProblemPhyXMini0559.boxVolumesInRestAndLaboratoryFrames}
\uses{def:physics:phyx-mini-0559:phyxminiproblems-problemphyxmini0559-volumeincubicmeters, def:physics:phyx-mini-0559:phyxminiproblems-problemphyxmini0559-relativisticrectangularboxsetup, def:physics:phyx-mini-0559:phyxminiproblems-problemphyxmini0559-matchesrectangularboxscenario, def:physics:phyx-mini-0559:phyxminiproblems-problemphyxmini0559-matchessuppliedboxfigure, def:physics:phyx-mini-0559:phyxminiproblems-problemphyxmini0559-matchesproblemreadouts, def:physics:phyx-mini-0559:phyxminiproblems-problemphyxmini0559-hasphysicalrelativisticboxparameters, def:physics:phyx-mini-0559:phyxminiproblems-problemphyxmini0559-satisfiesrectangularboxvolumegeometry, def:physics:phyx-mini-0559:phyxminiproblems-problemphyxmini0559-satisfiesspecialrelativisticvolumecontraction, def:physics:phyx-mini-0559:phyxminiproblems-problemphyxmini0559-agreeswhenroundedtoonedecimal, def:physics:phyx-mini-0559:phyxminiproblems-problemphyxmini0559-isnearestanswerchoice}
The assigned autoformalize agent should translate this physics problem into Lean declarations in the covered file, with theorem and lemma proof bodies written as `by sorry`.
\end{theorem}
\begin{proof}
This is an autoformalization task, not a proof task. Produce statements that can later be proved by the physics prover without weakening the physical model.
\end{proof}

% --- Archon declaration topology begin ---
\section{Declaration topology}
\begin{definition}[Length Quantity]
  \label{def:physics:phyx-mini-0559:phyxminiproblems-problemphyxmini0559-lengthquantity}
  \lean{PhyXMiniProblems.ProblemPhyXMini0559.LengthQuantity}
  A nonnegative, unit-independent physical length
\end{definition}
\begin{proof}
  This is the declared model definition of Length Quantity.
\end{proof}
\begin{definition}[Volume Quantity]
  \label{def:physics:phyx-mini-0559:phyxminiproblems-problemphyxmini0559-volumequantity}
  \lean{PhyXMiniProblems.ProblemPhyXMini0559.VolumeQuantity}
  A nonnegative physical volume carrying dimension L³
\end{definition}
\begin{proof}
  This is the declared model definition of Volume Quantity.
\end{proof}
\begin{definition}[length Readout]
  \label{def:physics:phyx-mini-0559:phyxminiproblems-problemphyxmini0559-lengthreadout}
  \lean{PhyXMiniProblems.ProblemPhyXMini0559.lengthReadout}
  \uses{def:physics:phyx-mini-0559:phyxminiproblems-problemphyxmini0559-lengthquantity}
  Read a physical length in a selected length unit
\end{definition}
\begin{proof}
  This is the declared model definition of length Readout.
\end{proof}
\begin{definition}[volume Readout]
  \label{def:physics:phyx-mini-0559:phyxminiproblems-problemphyxmini0559-volumereadout}
  \lean{PhyXMiniProblems.ProblemPhyXMini0559.volumeReadout}
  \uses{def:physics:phyx-mini-0559:phyxminiproblems-problemphyxmini0559-volumequantity}
  Read a physical volume in the cube of a selected length unit
\end{definition}
\begin{proof}
  This is the declared model definition of volume Readout.
\end{proof}
\begin{definition}[length In Meters]
  \label{def:physics:phyx-mini-0559:phyxminiproblems-problemphyxmini0559-lengthinmeters}
  \lean{PhyXMiniProblems.ProblemPhyXMini0559.lengthInMeters}
  \uses{def:physics:phyx-mini-0559:phyxminiproblems-problemphyxmini0559-lengthquantity, def:physics:phyx-mini-0559:phyxminiproblems-problemphyxmini0559-lengthreadout}
  Metre readout used for the three proper edge lengths
\end{definition}
\begin{proof}
  This is the declared model definition of length In Meters.
\end{proof}
\begin{definition}[volume In Cubic Meters]
  \label{def:physics:phyx-mini-0559:phyxminiproblems-problemphyxmini0559-volumeincubicmeters}
  \lean{PhyXMiniProblems.ProblemPhyXMini0559.volumeInCubicMeters}
  \uses{def:physics:phyx-mini-0559:phyxminiproblems-problemphyxmini0559-volumequantity, def:physics:phyx-mini-0559:phyxminiproblems-problemphyxmini0559-volumereadout}
  Cubic-metre readout used for both requested volumes and the choices
\end{definition}
\begin{proof}
  This is the declared model definition of volume In Cubic Meters.
\end{proof}
\begin{definition}[Inertial Frame Label]
  \label{def:physics:phyx-mini-0559:phyxminiproblems-problemphyxmini0559-inertialframelabel}
  \lean{PhyXMiniProblems.ProblemPhyXMini0559.InertialFrameLabel}
  The laboratory frame and the box's proper rest frame
\end{definition}
\begin{proof}
  This is the declared model definition of Inertial Frame Label.
\end{proof}
\begin{definition}[Proper Edge Label]
  \label{def:physics:phyx-mini-0559:phyxminiproblems-problemphyxmini0559-properedgelabel}
  \lean{PhyXMiniProblems.ProblemPhyXMini0559.ProperEdgeLabel}
  The three proper edge labels printed beside the box
\end{definition}
\begin{proof}
  This is the declared model definition of Proper Edge Label.
\end{proof}
\begin{definition}[Primed Coordinate Axis]
  \label{def:physics:phyx-mini-0559:phyxminiproblems-problemphyxmini0559-primedcoordinateaxis}
  \lean{PhyXMiniProblems.ProblemPhyXMini0559.PrimedCoordinateAxis}
  The primed Cartesian axes visible in the supplied raster
\end{definition}
\begin{proof}
  This is the declared model definition of Primed Coordinate Axis.
\end{proof}
\begin{definition}[Boost Direction]
  \label{def:physics:phyx-mini-0559:phyxminiproblems-problemphyxmini0559-boostdirection}
  \lean{PhyXMiniProblems.ProblemPhyXMini0559.BoostDirection}
  Direction selected for the standard relative boost
\end{definition}
\begin{proof}
  This is the declared model definition of Boost Direction.
\end{proof}
\begin{definition}[Volume Measurement Protocol]
  \label{def:physics:phyx-mini-0559:phyxminiproblems-problemphyxmini0559-volumemeasurementprotocol}
  \lean{PhyXMiniProblems.ProblemPhyXMini0559.VolumeMeasurementProtocol}
  Protocol used to assign a three-dimensional volume in an inertial frame
\end{definition}
\begin{proof}
  This is the declared model definition of Volume Measurement Protocol.
\end{proof}
\begin{definition}[Relativistic Box Figure]
  \label{def:physics:phyx-mini-0559:phyxminiproblems-problemphyxmini0559-relativisticboxfigure}
  \lean{PhyXMiniProblems.ProblemPhyXMini0559.RelativisticBoxFigure}
  \uses{def:physics:phyx-mini-0559:phyxminiproblems-problemphyxmini0559-properedgelabel, def:physics:phyx-mini-0559:phyxminiproblems-problemphyxmini0559-primedcoordinateaxis}
  Qualitative and scalar information transcribed from phyx\_data/test\_image/559.png. The angle is a dimensionless degree readout; the actual physical edge lengths and volumes are stored in the setup below
\end{definition}
\begin{proof}
  This is the declared model definition of Relativistic Box Figure.
\end{proof}
\begin{definition}[Relativistic Rectangular Box Setup]
  \label{def:physics:phyx-mini-0559:phyxminiproblems-problemphyxmini0559-relativisticrectangularboxsetup}
  \lean{PhyXMiniProblems.ProblemPhyXMini0559.RelativisticRectangularBoxSetup}
  \uses{def:physics:phyx-mini-0559:phyxminiproblems-problemphyxmini0559-lengthquantity, def:physics:phyx-mini-0559:phyxminiproblems-problemphyxmini0559-volumequantity, def:physics:phyx-mini-0559:phyxminiproblems-problemphyxmini0559-inertialframelabel, def:physics:phyx-mini-0559:phyxminiproblems-problemphyxmini0559-properedgelabel, def:physics:phyx-mini-0559:phyxminiproblems-problemphyxmini0559-boostdirection, def:physics:phyx-mini-0559:phyxminiproblems-problemphyxmini0559-volumemeasurementprotocol, def:physics:phyx-mini-0559:phyxminiproblems-problemphyxmini0559-relativisticboxfigure}
  The physical box and its independent observables. In particular, laboratoryVolume is not defined from the contraction formula or an answer choice; it is constrained only by the governing-law hypothesis below
\end{definition}
\begin{proof}
  This is the declared model definition of Relativistic Rectangular Box Setup.
\end{proof}
\begin{definition}[lorentz Factor]
  \label{def:physics:phyx-mini-0559:phyxminiproblems-problemphyxmini0559-lorentzfactor}
  \lean{PhyXMiniProblems.ProblemPhyXMini0559.lorentzFactor}
  \uses{def:physics:phyx-mini-0559:phyxminiproblems-problemphyxmini0559-relativisticrectangularboxsetup}
  Physlib's Lorentz factor for the stated dimensionless relative speed
\end{definition}
\begin{proof}
  This is the declared model definition of lorentz Factor.
\end{proof}
\begin{definition}[Matches Rectangular Box Scenario]
  \label{def:physics:phyx-mini-0559:phyxminiproblems-problemphyxmini0559-matchesrectangularboxscenario}
  \lean{PhyXMiniProblems.ProblemPhyXMini0559.MatchesRectangularBoxScenario}
  \uses{def:physics:phyx-mini-0559:phyxminiproblems-problemphyxmini0559-relativisticrectangularboxsetup}
  The box is at rest in S', while its volume in S is measured at one lab time
\end{definition}
\begin{proof}
  This is the declared model definition of Matches Rectangular Box Scenario.
\end{proof}
\begin{definition}[Matches Supplied Box Figure]
  \label{def:physics:phyx-mini-0559:phyxminiproblems-problemphyxmini0559-matchessuppliedboxfigure}
  \lean{PhyXMiniProblems.ProblemPhyXMini0559.MatchesSuppliedBoxFigure}
  \uses{def:physics:phyx-mini-0559:phyxminiproblems-problemphyxmini0559-relativisticrectangularboxsetup}
  Primary-image evidence: all three primed axes and edge labels are visible, the blue solid is rectangular, S' labels the primed frame, and the marked orientation relative to x' is 25°
\end{definition}
\begin{proof}
  This is the declared model definition of Matches Supplied Box Figure.
\end{proof}
\begin{definition}[Matches Problem Readouts]
  \label{def:physics:phyx-mini-0559:phyxminiproblems-problemphyxmini0559-matchesproblemreadouts}
  \lean{PhyXMiniProblems.ProblemPhyXMini0559.MatchesProblemReadouts}
  \uses{def:physics:phyx-mini-0559:phyxminiproblems-problemphyxmini0559-lengthinmeters, def:physics:phyx-mini-0559:phyxminiproblems-problemphyxmini0559-relativisticrectangularboxsetup}
  Numerical data supplied by the prose. These fields mention neither requested volume nor any answer choice
\end{definition}
\begin{proof}
  This is the declared model definition of Matches Problem Readouts.
\end{proof}
\begin{definition}[Has Physical Relativistic Box Parameters]
  \label{def:physics:phyx-mini-0559:phyxminiproblems-problemphyxmini0559-hasphysicalrelativisticboxparameters}
  \lean{PhyXMiniProblems.ProblemPhyXMini0559.HasPhysicalRelativisticBoxParameters}
  \uses{def:physics:phyx-mini-0559:phyxminiproblems-problemphyxmini0559-lengthinmeters, def:physics:phyx-mini-0559:phyxminiproblems-problemphyxmini0559-volumeincubicmeters, def:physics:phyx-mini-0559:phyxminiproblems-problemphyxmini0559-relativisticrectangularboxsetup}
  Positivity of the box observables and the physical subluminal regime
\end{definition}
\begin{proof}
  This is the declared model definition of Has Physical Relativistic Box Parameters.
\end{proof}
\begin{definition}[Satisfies Rectangular Box Volume Geometry]
  \label{def:physics:phyx-mini-0559:phyxminiproblems-problemphyxmini0559-satisfiesrectangularboxvolumegeometry}
  \lean{PhyXMiniProblems.ProblemPhyXMini0559.SatisfiesRectangularBoxVolumeGeometry}
  \uses{def:physics:phyx-mini-0559:phyxminiproblems-problemphyxmini0559-lengthreadout, def:physics:phyx-mini-0559:phyxminiproblems-problemphyxmini0559-volumereadout, def:physics:phyx-mini-0559:phyxminiproblems-problemphyxmini0559-relativisticrectangularboxsetup}
  Euclidean volume of a rectangular box is the product of its three mutually orthogonal proper edge lengths. Stating this in every unit keeps it a generic geometry law rather than a problem-specific numerical answer
\end{definition}
\begin{proof}
  This is the declared model definition of Satisfies Rectangular Box Volume Geometry.
\end{proof}
\begin{definition}[Satisfies Special Relativistic Volume Contraction]
  \label{def:physics:phyx-mini-0559:phyxminiproblems-problemphyxmini0559-satisfiesspecialrelativisticvolumecontraction}
  \lean{PhyXMiniProblems.ProblemPhyXMini0559.SatisfiesSpecialRelativisticVolumeContraction}
  \uses{def:physics:phyx-mini-0559:phyxminiproblems-problemphyxmini0559-volumereadout, def:physics:phyx-mini-0559:phyxminiproblems-problemphyxmini0559-relativisticrectangularboxsetup, def:physics:phyx-mini-0559:phyxminiproblems-problemphyxmini0559-lorentzfactor}
  For a volume obtained from one simultaneity slice in the laboratory frame, the determinant of the spatial Lorentz contraction gives V\_S = V\_S' / γ(β). This generic law is independent of the box's rest-frame orientation and contains no problem-specific volume value
\end{definition}
\begin{proof}
  This is the declared model definition of Satisfies Special Relativistic Volume Contraction.
\end{proof}
\begin{definition}[Answer Choice]
  \label{def:physics:phyx-mini-0559:phyxminiproblems-problemphyxmini0559-answerchoice}
  \lean{PhyXMiniProblems.ProblemPhyXMini0559.AnswerChoice}
  Labels of the four volume choices printed in the source problem
\end{definition}
\begin{proof}
  This is the declared model definition of Answer Choice.
\end{proof}
\begin{definition}[displayed Volume Cubic Meters]
  \label{def:physics:phyx-mini-0559:phyxminiproblems-problemphyxmini0559-displayedvolumecubicmeters}
  \lean{PhyXMiniProblems.ProblemPhyXMini0559.displayedVolumeCubicMeters}
  \uses{def:physics:phyx-mini-0559:phyxminiproblems-problemphyxmini0559-answerchoice}
  Displayed cubic-metre value beside each answer label
\end{definition}
\begin{proof}
  This is the declared model definition of displayed Volume Cubic Meters.
\end{proof}
\begin{definition}[recorded Dataset Answer]
  \label{def:physics:phyx-mini-0559:phyxminiproblems-problemphyxmini0559-recordeddatasetanswer}
  \lean{PhyXMiniProblems.ProblemPhyXMini0559.recordedDatasetAnswer}
  \uses{def:physics:phyx-mini-0559:phyxminiproblems-problemphyxmini0559-answerchoice}
  The answer label recorded by the source dataset
\end{definition}
\begin{proof}
  This is the declared model definition of recorded Dataset Answer.
\end{proof}
\begin{definition}[Agrees When Rounded To One Decimal]
  \label{def:physics:phyx-mini-0559:phyxminiproblems-problemphyxmini0559-agreeswhenroundedtoonedecimal}
  \lean{PhyXMiniProblems.ProblemPhyXMini0559.AgreesWhenRoundedToOneDecimal}
  \uses{def:physics:phyx-mini-0559:phyxminiproblems-problemphyxmini0559-answerchoice, def:physics:phyx-mini-0559:phyxminiproblems-problemphyxmini0559-displayedvolumecubicmeters}
  Agreement with a volume printed to one decimal place
\end{definition}
\begin{proof}
  This is the declared model definition of Agrees When Rounded To One Decimal.
\end{proof}
\begin{definition}[Is Nearest Answer Choice]
  \label{def:physics:phyx-mini-0559:phyxminiproblems-problemphyxmini0559-isnearestanswerchoice}
  \lean{PhyXMiniProblems.ProblemPhyXMini0559.IsNearestAnswerChoice}
  \uses{def:physics:phyx-mini-0559:phyxminiproblems-problemphyxmini0559-answerchoice, def:physics:phyx-mini-0559:phyxminiproblems-problemphyxmini0559-displayedvolumecubicmeters}
  A choice is at least as close to the computed lab volume as every alternative
\end{definition}
\begin{proof}
  This is the declared model definition of Is Nearest Answer Choice.
\end{proof}
% --- Archon declaration topology end ---
% --- Archon physics formalization source end ---
